\section{Comparative Statics: Rights, Costs and Resumption}
\label{sec:comparative}

The covered call values also price changes in the problem itself. In this section consider two per-arm problems $\mathcal P$ and $\mathcal P'$, each satisfying Assumption~\ref{assm:standing}, with the same arm path laws, so that each arm's readings are the same random variables under both. Say that $\mathcal P'$ \emph{dominates} $\mathcal P$ if $\cK_i\subseteq\cK_i'$ and
$c_i'(s)\le c_i(s)$ for every arm and own state: more readings may be kept and
no advance costs more. Corresponding states share the observed context, hence the own states $s_i$, and have incumbents $q\le q'$; the state of $\mathcal P'$ corresponding to $S$ is written $\bar S$. Primes on $W$, $Z$, $G$, $\Phi$, $\Yc$, $c$, $q$, $V$ and $U$ denote quantities of $\mathcal P'$ throughout this section, while $s'$ keeps its meaning as the successor own state.

\begin{lemma}[Pathwise monotonicity]\label{lem:pathwise}
If $\mathcal P'$ dominates $\mathcal P$, then at every own state $s$,
$W'(s)\ge W(s)$ and $Z'(s)\ge Z(s)$ on every remaining path, and $G'(s)\ge G(s)$.
\end{lemma}
\begin{proof}
Backward induction on $\nu(s)$; all three are $-\infty$ when $\nu(s)=0$. At $s$
with $\nu(s)\ge1$, the collectable reading satisfies $\Yc'\ge\Yc$ because
$\cK\subseteq\cK'$, and $Z'(s')\ge Z(s')$ by the hypothesis at $s'$, so
$W'(s)=\max\{\Yc',Z'(s')\}\ge W(s)$. By Lemma~\ref{lem:nested},
$\Phi'(s;r)=\E[(W'(s)-r)^+]-c'(s)\ge\E[(W(s)-r)^+]-c(s)=\Phi(s;r)$ for every
$r$, so the positivity set of $\Phi'$ contains that of $\Phi$ and
Proposition~\ref{prop:index} gives $G'(s)\ge G(s)$. Then
$Z'(s)=\min\{G'(s),W'(s)\}\ge\min\{G(s),W(s)\}=Z(s)$.
\end{proof}

\begin{proposition}[Exact comparative statics]\label{prop:comparative}
Let $\mathcal P'$ dominate $\mathcal P$. At corresponding states $S$ and $\bar S$, with $U=V^\star+K$ as in Section~\ref{sec:setting}, \[ U'(\bar S)-U(S)=\E\Bigl[\bigl(\max\{q',\max_iZ_i'(s_i)\}-\max\{q,\max_iZ_i(s_i)\}\bigr)^+\Bigr]\ge0,
\]
the expectation coupling $Z_i$ and $Z_i'$ on arm $i$'s remaining path under
$\mathsf P_S$. At a common initial state the value of the additional rights or
lower costs is $\E[(\max_iZ_i'-\max\{q,\max_iZ_i\})^+]$, which vanishes if and
only if $\max_iZ_i'\le\max\{q,\max_iZ_i\}$ almost surely.
\end{proposition}
\begin{proof}
Theorem~\ref{thm:index} applies to both problems on the common product space, giving $U'(\bar S)-U(S)=\E[\max\{q',\max_iZ_i'\}-\max\{q,\max_iZ_i\}]$. By
Lemma~\ref{lem:pathwise} and $q'\ge q$ the integrand is nonnegative on every
path, so it equals its positive part.
\end{proof}

The identity follows from pathwise monotonicity and the value representation. Its use is
to price additional collection rights or lower costs from coupled arm-wise covered-call
values, without solving either joint Bellman problem. It covers enlarging a collectable
set, removing a history-dependent surcharge, and lowering any own-history cost.
For one-shot boxes, adding a box or lowering a cost raises the stopping frontier
\citep[Proposition~2]{bao2026interdependent}. The identity does not cover changes to the
readings' law, such as adding an observation stage or changing signal accuracy.
The next result prices the opposite error: planning
as if fewer readings could be kept.

\begin{proposition}[Planning with a smaller collectable set]\label{prop:fewer-rights}
Let $\mathcal P'$ dominate $\mathcal P$ with the same costs, so that only the collectable sets differ. Let $\pi_{\cK}$ be the policy for $\mathcal P'$ that stops when $\max_iG_i(s_i)\le q'$ and otherwise advances an arm maximizing the $\mathcal P$-index $G_i(s_i)$, where $q'$ is the incumbent of $\mathcal P'$. Then at corresponding states \[ V^\star(S)\le V_{\pi_{\cK}}(\bar S)\le V^{\star\prime}(\bar S). \] \end{proposition} Its loss against the optimum of $\mathcal P'$ is therefore at most the value of the additional rights; both ends are attained in the study of Appendix~\ref{app:comparative-results}.

Appendix~\ref{app:comparative} proves this and the next result. Finally, the
loss decomposition prices the freedom to return to a paused arm. Fix a starting state
$S_0$. A policy \emph{resumes} if it returns to an arm that it has advanced since $S_0$
after advancing a different arm; a \emph{no-return} policy never resumes, and $V_{\rm NR}(S_0)$
is its optimal value. This differs from search without recall, or committed search
\citep[Section~5]{gupta2019markovian}, which forbids returning to an earlier prize:
a no-return policy retains every collectable reading.

On a path where the index resumes, a no-return policy must depart from its
trajectory by that advance. We lower-bound the resulting loss by the smallest
expected Bellman penalty at this first departure; later losses are nonnegative.

To preserve return status in the bound, augment a descendant history with
$\widehat S=(S,\ell,D)$, where $\ell$ is the last arm advanced since $S_0$
($\ell=0$ initially) and $D$ is the set of arms advanced and subsequently left since
$S_0$. An advance of $i$ sets $\ell^+=i$ and
$D^+=D\cup\{\ell\}$ if $\ell\notin\{0,i\}$, and $D^+=D$ otherwise.
It is a resumption exactly when $i\in D$ and $i\ne\ell$. This bookkeeping is retained along the
whole run, including when the last arm is exhausted.
For the index policy with its fixed tie order, let $\delta(S)$ be the smallest
loss from Theorem~\ref{thm:loss} of an action other than $\pi_{\rm IX}(S)$,
with $\min\varnothing=+\infty$. Define $\widehat\rho(\widehat S)$ by finite backward
recursion: it is zero if the index policy has zero probability of a future resumption
from this augmented history, equals $\delta(S)$ if its next action resumes, and otherwise
\[
 \widehat\rho(\widehat S)
 =\min\{\delta(S),\E[\widehat\rho(\widehat S_1)\mid\widehat S]\},
 \qquad \rho(S_0)=\widehat\rho(S_0,0,\varnothing),
\]
where $\widehat S_1$ is the successor under the index policy with the bookkeeping above.
The minimum compares departing now, at penalty $\delta(S)$, with following the
index one more step and facing a later departure. At the first resumption,
departure is compulsory. This first-departure calculation gives a lower bound
on the total resumption gap.

\begin{proposition}[Value of resumption]\label{prop:resumption}
Under the hypotheses of Theorem~\ref{thm:index}:
\begin{enumerate}
\item[(i)] $V^\star(S_0)-V_{\rm NR}(S_0)\ge\rho(S_0)\ge0$, and $V_{\rm NR}(S_0)=V^\star(S_0)$ when $\pi_{\rm IX}$ almost surely never resumes from $S_0$.
\item[(ii)] If at every state at which $\pi_{\rm IX}$ advances from $S_0$, before or at its first resumption, the index-maximal arm is unique, then $\rho(S_0)>0$ exactly when $\pi_{\rm IX}$ resumes from $S_0$
 with positive probability.
\item[(iii)] If every reading is a deterministic function of the arm's own
 state, $\pi_{\rm IX}$ never resumes under any tie order, so $V_{\rm NR}=V^\star$.
\end{enumerate}
\end{proposition}

With deterministic readings an advance either raises the arm's index by its cost
or lifts the incumbent above its previous index, which rules out resumption.
Stochastic observations can make resumption valuable by resolving uncertainty in
collectable rewards, even without changing any unobserved reward's law.
Example~\ref{ex:resumption-reward} in Appendix~\ref{app:comparative} gives a
two-arm construction and calculates the values with and without resumption.

Section~\ref{sec:validation} and Appendix~\ref{app:comparative-results}
evaluate all three results exactly on the staging grid.
